\section{Theory}
\label{sec:theory}

This section states the results and the intuition; every proof is in Appendix~\ref{app:proofs}, which also carries the
map from the numbering here to the numbering of the technical notes.

\subsection{An exact test, at any sample size}

\begin{theorem}[Exactness of the flip test]
\label{thm:exact}
Let $Z_1,\dots,Z_K$ be data elements that are exchangeable under $H_0$, and let $T$ be any statistic. Draw $M-1$ i.i.d.\
uniform group elements $\sigma_m\in\symgrp$ (equivalently, i.i.d.\ signs on the antisymmetric component) and let
$p=\frac{1}{M}\bigl(1+\#\{m: T(\rho(\sigma_m)Z)\ge T(Z)\}\bigr)$. Then under $H_0$, $\mathbb{P}(p\le\alpha)\le\alpha$
for every $K$ and every $M$, with equality when $\alpha M$ is an integer.
\end{theorem}
The mechanism is that under $H_0$ the observed configuration and its group images are equally likely, so the rank of the
observed statistic among its own orbit is uniform. Nothing is estimated, nothing is fitted, and no asymptotics are
invoked: the level holds at $K=10$ as it does at $K=1000$. The price is that the test only sees the group --- which is
exactly what the next theorem characterises.

\subsection{What is visible: a two-layer characterisation}

\begin{theorem}[Two-layer detectability]
\label{thm:twolayer}
\emph{(Layer I, law level.)} If the fault pattern is fixed by $\symgrp$ (a bilaterally symmetric degradation) \emph{and}
the faulted closed loop retains a unique symmetric attractor, then $\push{\rho}P_\varphi=P_\varphi$ and \emph{every}
level-$\alpha$ invariance test has power $\le\alpha$, whatever the magnitude. Conversely, if
$\push{\rho}P_\varphi\ne P_\varphi$, a statistic consistent for the two-sample problem between $\{Z_k\}$ and
$\{\rho Z_k\}$ --- the energy-distance class --- has power $\to 1$.
\emph{(Layer II, mean level.)} For the paired-energy statistic the noncentrality after $K$ elements is
$K\|\Pi^-\mu(\varphi)\|^2_\Sigma$, and the test is consistent \emph{iff} $\Pi^-\mu(\varphi)\ne 0$ --- a strict subset of
the laws Layer~I covers.
\end{theorem}

Three consequences drive the system design. First, blindness to bilaterally symmetric faults is structural, not a
weakness of a particular statistic, so a second, magnitude channel reading $\mathcal{W}^+$ is \emph{necessary}, and it
cannot inherit the same nuisance immunity. Second, the two layers differ exactly on zero-mean law changes: a one-sided
variance inflation has $\Pi^-\mu=0$ but breaks the law, so the paired statistic is blind while the energy-distance one
is consistent --- deployment runs both. Third, the blindness is robust in amplitude: sweeping a bilateral fault down to
$30\,\%$ of nominal torque left the test at its level and the realised gait's chirality at its nominal floor, so the
escape route is a genuine symmetry-breaking bifurcation of the gait, not a large enough fault.

\begin{proposition}[Power and sample size]
\label{prop:power}
Under a Gaussian model for the differenced scores, the power of the whitened quadratic flip statistic at level $\alpha$
is at least $1-\exp\bigl(-(d+\lambda-c_\alpha)_+^2/(4(d+2\lambda))\bigr)$ with $\lambda=K\|\Pi^-\mu\|^2_{\Sigma}$, which
inverts to an explicit sample-size rule $K\ge\lambda^*(d,\alpha,\beta)/\|\Pi^-\mu\|^2_{\Sigma}$.
\end{proposition}

\subsection{Reading it every cycle}

\begin{proposition}[Anytime-valid sequential layer]
\label{prop:seq}
Let $p_1,p_2,\dots$ be the window $p$-values of Thm.~\ref{thm:exact} and $e_t=\tfrac12 p_t^{-1/2}$. Then
$E_t=\prod_{s\le t}e_s$ is a nonnegative supermartingale with $E_0=1$, and by Ville's inequality the rule
$E_t\ge 1/\alpha$ has false-alarm probability at most $\alpha$ over an unbounded horizon.
\end{proposition}
This is what makes the test deployable: it may be inspected after every cycle, and stopping when it looks interesting
does not invalidate it.

\subsection{Residual channels and equivariance}

Running the same test on the residual of a nominal dynamics model, $r=y-\hat f(x,u)$, both sharpens it and endangers it.

\begin{proposition}[Residual inheritance and its contamination]
\label{prop:residual}
If the nominal model $\hat f$ is $\symgrp$-equivariant, the residual element inherits $H_0$ exactly, and the flip test
on it retains its level while gaining power against faults that act through the dynamics. If $\hat f$ is not
equivariant, its equivariance defect enters the residual as a fixed asymmetric offset and inflates the level.
\end{proposition}
The effect is large enough to matter in practice, not merely in principle: on the welded-leg world a plain learned model
drives the residual test's size to $1.00$ while its equivariant twin sits at $0.00$ (\S\ref{sec:simulation}). Mirror
weight sharing is what makes a Deep Lagrangian Network equivariant by construction, and the equivariance defect is then
exactly zero rather than small.

\subsection{Rolling contact keeps the invariant filter}

\begin{proposition}[Rolling contact is group-affine]
\label{prop:rolling}
For a wheel in contact with rolling constraint $\dot d_i=R\,u_i$, $u_i$ a measured input, the contact-aided invariant
filter's propagation remains group-affine, and the right-invariant error dynamics of the contact block is
\emph{unchanged}: $A[d_i,R]=0$, exactly as for a world-fixed foot. Rolling alters only the mean propagation
$d_i\mapsto d_i+Ru_i\Delta t$ and the contact process noise, which becomes anisotropic in the wheel frame.
\end{proposition}
The plain-EKF Jacobian for a moving world point is $-R(u_i)_\times$, and carrying it into an invariant filter is a
natural and wrong step that makes the error dynamics state-dependent. The correct statement admits wheeled-legged
robots to the same estimator, and hence the same gating front-end, as point-foot robots. The signature of a fault in
such a filter is the observable projection of its direction transported by the adjoint, which is why a bilaterally
symmetric slip is invisible to the invariance channel but visible to the estimator channel --- the property that turns
the pair into a slip \emph{classifier} (\S\ref{sec:method}).
